We have now discussed the foundations of LLMs and agents. We have seen how language and tokenization works, we have seen how LLMs are trained and finetuned, and we have seen how advanced reasoning models reason and self-reflect. We have looked at different learning modalities, from language, vision, and action.  We have also looked at the intricacies of metalearning. 

We have now arrived at the heart of agentic LLM, where we will review the different interaction styles of LLM  agents.  This part consists of one content chapter and one practical chapter. First we will review the five agentic interaction styles that we have seen so far, and we discuss the importance of memory. Then, in the practical chapter, we will discuss a  hands-on case study where many of the techniques of this part come together. You will be invited to create a self-reflective LLM agent that can play typical AI-games such as TextWorld and NetHack.
